% ============================================================
%  9. BIBLIOGRAFÍA Y ANEXOS   (guía, 2.9 y 3.6)
%  Bibliografía con biblatex + biber (Fase 4): estilo APA (2.6),
%  alfabética, sin numerar, sangría francesa, todo generado por
%  \printbibliography a partir de referencias.bib. Solo aparecen
%  las obras efectivamente citadas en el texto (\parencite en
%  secciones/03_problema.tex y 07_marco_teorico.tex); las demás
%  entradas de referencias.bib no citadas en el cuerpo quedan
%  fuera de la lista impresa, tal como exige "solo obras realmente
%  citadas" -- comportamiento nativo de biblatex, sin \nocite.
% ============================================================
\chapter{Bibliografía y Anexos}
\label{cap:bibliografia}

\section{Bibliografía}
\label{sec:bibliografia}

\printbibliography[heading=none]
% Nota interna sobre verificación de referencias y checklist de cierre:
% movidas a docs/notas_internas.md (defecto #9 — no van dentro del PDF).

\clearpage
\section{Anexos}
\label{sec:anexos}

Los anexos incorporan la evidencia técnica que respalda la monografía sin recargar el
cuerpo principal del documento.

\subsection{Anexo A. Diccionario de datos completo}
\label{anexo:diccionario-datos}

Ver tabla~\ref{tab:diccionario-datos} (sección~\ref{sec:diccionario-datos}) para el
diccionario completo del conjunto consolidado, incluyendo los campos de trazabilidad
(\var{source_file}, \var{source_batch}) añadidos durante la consolidación.

\subsection{Anexo B. Trazabilidad de notebooks y evidencias}
\label{anexo:trazabilidad-scripts}

Ver tabla~\ref{tab:trazabilidad-oe} (sección~\ref{sec:enlace-objetivos}) para la
correspondencia entre los objetivos específicos, los seis notebooks del pipeline y la
evidencia que produce cada uno. Estructura del repositorio del proyecto de datos:
\texttt{data/raw} (85 CSV de consultas, el feed GTFS y el archivo de Kontur, solo
lectura), \texttt{data/interim} (consultas limpias y tablas intermedias en formato Parquet),
\texttt{data/processed} (tabla de modelado por zona), \texttt{src} (funciones que llaman
los notebooks), \texttt{notebooks} (\var{01_eda} a \var{06_propuesta}) y
\texttt{resultados/<fase>} (tablas, figuras y \var{metricas.json} de cada fase).
